\subsection{Verification Scope}
An analytical lower bound certifies the T2 placement for the local RMS objective. Quadrature and Monte Carlo probes check the wrapped coverage crossing in Problem 3. The branch-and-bound result covers the 585-station discrete pool and the offset-corner graph. A separate tangent-arc comparison measures the waypoint error on the single audited nominal detour with both endpoints fixed.

\subsection{Limitations}
\textbf{Motion and separation.} The constant-speed disk model omits acceleration and curvature constraints. Candidate-route and executed-path audits check building clearance and field limits. Simultaneous paired stations in the nominal plan are at least 38.43~m apart. Dynamic vehicle-to-vehicle separation lies outside the dispatch constraints.

\textbf{Uncertainty and target motion.} The first-order RMS is a local measure. It cannot describe rare nonlinear intersections when noisy bearings are nearly collinear. The pointing model uses isotropic odometry at a fixed target range. Extending the results to moving targets or anisotropic odometry would require new tracking and measurement models.

\textbf{Scenarios and weights.} The 115 distinct disturbance cases cover specified faults, outage clocks, and spatial alert variations in a deterministic stress set. Completion fractions summarize those cases. The sensitivity index weights observation quality, dispatch latency, admissible distance, and recovery performance. Changing those weights changes some parameter rankings (Figure~\ref{fig:sens}d).
